\documentclass[11pt]{article}
\usepackage[margin=1in]{geometry}
\usepackage{amsmath,amssymb,amsthm,mathtools}
\usepackage{microtype}
\usepackage[hidelinks]{hyperref}

\newcommand{\M}{\mathcal M}
\newcommand{\N}{\mathcal N}
\newcommand{\X}{\mathfrak X}
\newcommand{\cL}{\mathcal L}
\newcommand{\id}{\mathrm{id}}
\newcommand{\modcl}{\operatorname{mod}}

\title{Surjective Poisson Submersions, Relative Poisson Flow, and Modular Dynamics}
\author{}
\date{}

\begin{document}
\maketitle
\vspace{-1.5em}

\section*{1. Poisson submersion, gauge symmetry, and Poisson cohomology}

Let
\[
p:(M,\pi_M)\longrightarrow (B,\pi_B)
\]
be a surjective Poisson submersion, so that
\[
\{f,g\}_B\circ p=\{p^*f,p^*g\}_M,
\qquad f,g\in C^\infty(B).
\]
Equivalently, the pullback identifies the reduced observables with a Poisson
subalgebra
\[
p^*C^\infty(B)\subset C^\infty(M).
\]
This is the classical coarse-graining picture that will be compared with an
inclusion of von Neumann algebras
\[
\N\subset\M.
\]

The vertical bundle
\[
V=\ker(dp)\subset TM
\]
contains precisely the infinitesimal motions invisible to the reduced
observables:
\[
X\in\Gamma(V)
\quad\Longrightarrow\quad
X(p^*f)=0,
\qquad f\in C^\infty(B).
\]
If $B=M/G$ for a free and proper Poisson action of a Lie group $G$, then $V$
is tangent to the $G$-orbits.  Hence vertical motion is gauge motion.

The natural cohomological framework is Poisson cohomology.  For a Poisson
tensor $\pi$ define
\[
d_\pi=[\pi,\cdot]_{\mathrm S},
\]
where $[\cdot,\cdot]_{\mathrm S}$ is the Schouten bracket.  Since
\[
[\pi,\pi]_{\mathrm S}=0,
\]
we have $d_\pi^2=0$, and
\[
H_\pi^k(M)
=
\frac{\ker(d_\pi:\X^k(M)\to\X^{k+1}(M))}
{\operatorname{im}(d_\pi:\X^{k-1}(M)\to\X^k(M))}.
\]
In degree one,
\[
H_\pi^1(M)
=
\frac{\{X\in\X(M):\cL_X\pi=0\}}
{\{X_f=\pi^\sharp(df):f\in C^\infty(M)\}}.
\]
Thus $H_\pi^1(M)$ measures Poisson symmetries modulo Hamiltonian ones.  This is
the correct classical receptacle for ``outer'' relative dynamics: Hamiltonian
vector fields are the Poisson analogue of inner generators, while a nonzero
class in $H_\pi^1$ represents a Poisson symmetry that cannot be generated by a
single Hamiltonian observable.

There is a parallel role for $H_\pi^2$.  If
\[
\pi_\varepsilon=\pi+\varepsilon\dot\pi+O(\varepsilon^2)
\]
is a first-order deformation of the Poisson structure, then
\[
[\pi,\dot\pi]_{\mathrm S}=0.
\]
Deformations of the form
\[
\dot\pi=\cL_X\pi=-d_\pi X
\]
are infinitesimal changes of coordinates.  Hence $H_\pi^2$ measures genuine
first-order deformations of the Poisson structure, whereas $H_\pi^1$ measures
genuine infinitesimal Poisson automorphisms.

Now let $X_0,X_1$ be Poisson vector fields on $M$, with flows
$\Phi_t^0,\Phi_t^1$.  Their relative flow is
\[
C_t=(\Phi_t^0)^{-1}\circ\Phi_t^1.
\]
In general $C_t$ is not a one-parameter group.  Its time-dependent generator is
\[
Y_t
=
\frac{d}{dt}C_t\circ C_t^{-1}
=
(\Phi_t^0)^{-1}_*(X_1-X_0).
\]
In particular,
\[
C_t=\id+t(X_1-X_0)+O(t^2).
\]
If both flows preserve the same Poisson tensor, then
\[
\cL_{Y_t}\pi=0,
\]
so that
\[
[Y_t]\in H_\pi^1(M).
\]
When $Y_t$ is Hamiltonian, the relative dynamics is Poisson-inner; when
$[Y_t]\neq0$, the relative dynamics contains genuinely outer information.

The gauge interpretation now becomes precise.  Suppose $X_0$ and $X_1$ are
two lifts of the same reduced vector field $X_B$:
\[
dp(X_0)=X_B\circ p,
\qquad
dp(X_1)=X_B\circ p.
\]
Then
\[
dp(X_1-X_0)=0,
\]
hence
\[
X_1-X_0\in\Gamma(\ker dp).
\]
Therefore the infinitesimal relative dynamics is vertical and consequently
gauge.  In general, however, the relative generator need not be vertical.
Schematically one may write
\[
Y_t=Y_t^{\mathrm{hor}}+Y_t^{\mathrm{vert}},
\]
where $Y_t^{\mathrm{vert}}$ is gauge motion and $Y_t^{\mathrm{hor}}$ measures a
genuine mismatch between the reduced dynamics.

\section*{2. Modular flow and the Borchers relative product}

Let
\[
\N\subset\M
\]
be von Neumann algebras with a common cyclic and separating vector $\Omega$.
Their modular groups are
\[
\sigma_t^\M(A)
=
\Delta_\M^{it}A\Delta_\M^{-it},
\qquad
\sigma_t^\N(B)
=
\Delta_\N^{it}B\Delta_\N^{-it}.
\]
The natural relative object is the product
\[
\Gamma_{\N,\M}(t)
=
\Delta_\N^{it}\Delta_\M^{-it},
\]
up to the choice of ordering convention.  This compares two canonical modular
flows determined by the same reference vector but by different observable
algebras.

The essential point is that $\Gamma_{\N,\M}(t)$ is generally not obtained by
simply exponentiating the difference of two generators.  Writing formally
\[
\Delta_\M^{it}=e^{itK_\M},
\qquad
\Delta_\N^{it}=e^{itK_\N},
\]
one has
\[
\Gamma_{\N,\M}(t)
=
e^{itK_\N}e^{-itK_\M}
=
1+it(K_\N-K_\M)+O(t^2),
\]
whereas at finite time the noncommutativity of $K_\N$ and $K_\M$ matters.

This is exactly the same structural phenomenon as for the classical relative
flow
\[
C_t=(\Phi_t^0)^{-1}\Phi_t^1.
\]
Infinitesimally,
\[
C_t
=
\id+t(X_1-X_0)+O(t^2),
\]
so the first-order correspondence is
\[
K_\N-K_\M
\qquad\longleftrightarrow\qquad
X_1-X_0.
\]
At finite time both sides are naturally relative cocycles rather than flows
generated by a simple difference.

The comparison can therefore be summarized as
\[
\begin{aligned}
p^*C^\infty(B)\subset C^\infty(M)
&\quad\longleftrightarrow\quad
\N\subset\M,\\[2mm]
\Phi_t
&\quad\longleftrightarrow\quad
\Delta^{it},\\[2mm]
(\Phi_t^0)^{-1}\Phi_t^1
&\quad\longleftrightarrow\quad
\Delta_\N^{it}\Delta_\M^{-it},\\[2mm]
X_1-X_0
&\quad\longleftrightarrow\quad
K_\N-K_\M,\\[2mm]
[Y_t]\in H_\pi^1(M)
&\quad\longleftrightarrow\quad
\text{outer relative dynamical information},\\[2mm]
Y_t\in\ker(dp)
&\quad\longleftrightarrow\quad
\text{relative motion invisible to coarse observables}.
\end{aligned}
\]

The last line is especially important.  It explains why the gauge analogy is
only partial.  A relative modular product such as
$\Delta_\N^{it}\Delta_\M^{-it}$ need not represent a redundancy.  It may
contain genuine information about the inclusion $\N\subset\M$.  Accordingly,
the classical relative flow should not be forced to be vertical.  Its vertical
component models gauge redundancy; its nonvertical component models genuine
relative dynamics between two levels of coarse-graining.

\section*{3. The Poisson modular class as the intrinsic bridge}

The preceding construction used arbitrary Poisson vector fields.  There is,
however, an intrinsic Poisson-geometric object whose name and role already
suggest the correct connection with Tomita--Takesaki theory.

Choose a volume form $\mu$ on $(M,\pi)$.  For a Hamiltonian vector field $X_f$,
define
\[
X_\mu(f)
=
\operatorname{div}_\mu(X_f).
\]
The resulting vector field $X_\mu$ is a Poisson vector field:
\[
\cL_{X_\mu}\pi=0.
\]
If the density is changed according to
\[
\mu'=e^g\mu,
\]
then, up to the sign convention in the definition of $X_g$,
\[
X_{\mu'}=X_\mu-X_g.
\]
Thus changing the reference density changes the modular vector field only by a
Hamiltonian vector field.  Consequently
\[
\modcl(M,\pi)
=
[X_\mu]\in H_\pi^1(M)
\]
is independent of $\mu$.  This is the Poisson modular class.

The analogy with von Neumann modular theory is now much sharper.  The Poisson
modular vector field measures the failure of a reference density to be
invariant under Hamiltonian flows.  The Tomita--Takesaki modular group measures
the failure of a state or weight to behave tracially.  On the Poisson side,
changing the density changes the modular vector field by a Hamiltonian term;
on the operator-algebraic side, changing the reference weight leads naturally
to relative modular cocycles.  In both cases the most intrinsic information is
therefore relative and is defined only modulo an inner ambiguity.

For the Poisson submersion
\[
p:(M,\pi_M)\to(B,\pi_B),
\]
choose densities $\mu_M$ and $\mu_B$, with modular vector fields
$X_{\mu_M}$ and $X_{\mu_B}$.  If $\widetilde X_{\mu_B}$ is a lift of the
downstairs modular vector field, define the relative modular vector field
\[
R_{\mu_M,\mu_B}
=
X_{\mu_M}-\widetilde X_{\mu_B}.
\]
Changing the lift changes $R_{\mu_M,\mu_B}$ by a vertical vector field, while
changing the densities changes it by Hamiltonian terms.  Thus the genuinely
invariant object is naturally a relative cohomology class rather than a single
preferred vector field.

If
\[
dp(X_{\mu_M})
=
X_{\mu_B}\circ p,
\]
then
\[
dp(R_{\mu_M,\mu_B})=0,
\]
so the relative modular vector field is vertical.  In a quotient realization
$B=M/G$, its flow is then gauge motion.  If instead
\[
dp(X_{\mu_M})
-
X_{\mu_B}\circ p
\neq0,
\]
the relative modular defect has a component visible on the reduced phase
space.  This is the classical counterpart of the fact that a Borchers relative
modular product can encode genuine information about an inclusion rather than
merely an internal redundancy.

The resulting conceptual picture is therefore
\[
\boxed{\text{surjective Poisson submersion}}
\quad\longleftrightarrow\quad
\boxed{\text{inclusion of observable algebras}},
\]
\[
\boxed{\text{relative Poisson cocycle}}
\quad\longleftrightarrow\quad
\boxed{\text{relative modular/Borchers cocycle}},
\]
and
\[
\boxed{\text{vertical part of the relative cocycle}}
\quad\longleftrightarrow\quad
\boxed{\text{gauge redundancy}}.
\]
Poisson cohomology then separates Hamiltonian, hence inner, motion from genuinely
outer relative dynamics, while the Poisson modular class provides the intrinsic
classical modular datum.  This makes the surjective Poisson-submersion picture
a natural classical model for relative modular dynamics under coarse-graining.


\section*{4. Dynamical information encoded by the relative cocycle}

The analogy with the Borchers characteristic function becomes more informative
when one asks what can be recovered from the relative Poisson cocycle.  Let
\[
C_t=(\Phi_t^0)^{-1}\circ\Phi_t^1,
\]
where $\Phi_t^0,\Phi_t^1$ are Poisson flows generated by $X_0,X_1$.  Its
instantaneous generator is
\[
Y_t=\dot C_t\circ C_t^{-1}
=(\Phi_t^0)^{-1}_*(X_1-X_0),
\qquad
Y_0=X_1-X_0.
\]
Thus the cocycle determines the relative infinitesimal generator, although it
does not in general determine $X_0$ and $X_1$ separately.

If both dynamics are Hamiltonian for the same Poisson tensor,
\[
X_0=X_{H_0},\qquad X_1=X_{H_1},
\]
then
\[
X_1-X_0=X_{H_1-H_0}.
\]
Hence the infinitesimal cocycle determines
\[
V:=H_1-H_0
\]
up to a Casimir,
\[
V\sim V+c,\qquad \pi^\sharp(dc)=0.
\]
On a connected symplectic manifold the ambiguity is only an additive constant.
The cocycle therefore determines the relative Hamiltonian modulo precisely the
functions invisible to the Poisson bracket.

At finite time the cocycle contains more information.  Unless the two flows
commute,
\[
C_t\neq\exp\bigl(t(X_1-X_0)\bigr).
\]
Instead,
\[
Y_t=(\Phi_t^0)^{-1}_*X_V.
\]
Equivalently, up to the convention for pullback by the Hamiltonian flow,
\[
Y_t=X_{V_t},
\qquad
V_t=V\circ\Phi_t^0.
\]
With the Liouville operator
\[
L_{H_0}f=\{H_0,f\},
\]
this may be written schematically as
\[
V_t=e^{-tL_{H_0}}V,
\]
with the sign determined by the flow convention.  Thus $C_t$ is naturally an
\emph{interaction-picture cocycle}: the perturbation $V=H_1-H_0$ is transported
by the reference $H_0$-dynamics before generating the relative evolution.

Expanding near $t=0$,
\[
Y_t
=
X_1-X_0
-t[X_0,X_1-X_0]
+\frac{t^2}{2}[X_0,[X_0,X_1-X_0]]
+\cdots.
\]
Since
\[
[X_f,X_g]=X_{\{f,g\}},
\]
successive derivatives probe
\[
V,\qquad
\{H_0,V\},\qquad
\{H_0,\{H_0,V\}\},\qquad\ldots
\]
modulo Casimirs.  Hence the finite-time cocycle sees the dynamical orbit of the
perturbation under the Liouville operator, rather than only the static
difference $H_1-H_0$.

This is parallel to
\[
\Gamma_{\N,\M}(t)
=e^{itK_\N}e^{-itK_\M}.
\]
Its first-order term contains $K_\N-K_\M$, while higher orders depend on
commutators with the reference modular generator.  The refined correspondence
is therefore
\[
\begin{aligned}
H_0&\longleftrightarrow K_\M,\\
V=H_1-H_0&\longleftrightarrow K_\N-K_\M,\\
e^{-tL_{H_0}}V
&\longleftrightarrow
\text{modularly transported relative generator},\\
(\Phi_t^0)^{-1}\Phi_t^1
&\longleftrightarrow
\Delta_\N^{it}\Delta_\M^{-it}.
\end{aligned}
\]
In this sense the Borchers characteristic function is more naturally compared
with an interaction cocycle than with an ordinary Hamiltonian flow.

Several qualitative facts can be inferred directly.  If
\[
C_t=\id\quad\text{for all }t,
\]
then $X_0=X_1$.  The class
\[
[Y_0]\in H_\pi^1(M)
\]
tests whether the infinitesimal relative dynamics is Hamiltonian: vanishing
means Poisson-inner relative motion, while a nonzero class detects an outer
Poisson component.  For a Poisson submersion $p:M\to B$, the condition
\[
dp(Y_t)=0
\]
means that the relative cocycle acts entirely along the fibres and therefore
\[
C_t^*(p^*f)=p^*f,\qquad f\in C^\infty(B).
\]
Thus one obtains
\[
\begin{aligned}
Y_t=0
&\quad\Longrightarrow\quad \text{identical dynamics},\\
Y_t\neq0,\quad dp(Y_t)=0
&\quad\Longrightarrow\quad
\text{different lifts of the same reduced dynamics},\\
dp(Y_t)\neq0
&\quad\Longrightarrow\quad
\text{different dynamics visible after reduction}.
\end{aligned}
\]

The group property also contains dynamical information.  If
$[X_0,X_1]=0$, then
\[
C_t=\exp\bigl(t(X_1-X_0)\bigr),
\qquad C_{t+s}=C_tC_s.
\]
Failure to reduce to such an ordinary one-parameter group therefore records
nontrivial interaction between the background dynamics and the perturbation,
controlled in the Hamiltonian case by Poisson brackets of $H_0$ with $V$.

This suggests an ergodic interpretation.  The natural object is not only $V$
but its orbit
\[
\{e^{-tL_{H_0}}V:t\in\mathbb R\}.
\]
For an invariant measure $\mu$, correlations such as
\[
\int_M V\,(V\circ\Phi_t^0)\,d\mu
\]
measure how the perturbation generating the relative cocycle is mixed by the
background dynamics.  Decay of these correlations, spectral properties of
$L_{H_0}$ on the cyclic subspace generated by $V$, and the behavior of iterated
Poisson brackets are natural dynamical invariants attached to the relative
cocycle.

There is nevertheless a reconstruction limitation: $C_t$ alone does not
generally determine the two Hamiltonian systems separately.  If one reference
flow is known, however, the other is recovered from
\[
\Phi_t^1=\Phi_t^0\circ C_t.
\]
Thus the cocycle should be regarded as a relative dynamical invariant, just as
the Borchers characteristic function compares two modular structures rather
than constituting a complete invariant of either one individually.

Formally, the strengthened analogy is
\[
\boxed{
\Delta_\N^{it}\Delta_\M^{-it}
\quad\longleftrightarrow\quad
\mathcal T\exp\left(\int_0^t X_{V_s}\,ds\right),
\qquad
V_s=e^{-sL_{H_0}}(H_1-H_0),
}
\]
where $\mathcal T$ denotes time ordering.  The Poisson cocycle therefore
encodes the perturbation in the interaction picture: its first derivative gives
the relative Hamiltonian vector field, its higher derivatives probe iterated
Poisson brackets with the background Hamiltonian, its Poisson cohomology class
distinguishes inner from outer relative dynamics, and its projection under $p$
distinguishes gauge motion from physically visible relative evolution.

\end{document}
